% ==============================================================================
% SECTION 2: RELATED WORK & BASELINE POSITIONING
% ==============================================================================

\section{Related Work \& Baseline Positioning}
\label{sec:related_work}

\subsection{Sequential Recommender Systems}
Sequential recommendation models user preference dynamics from chronological interaction histories $\mathcal{S}_u = [v_1, v_2, \dots, v_t]$. Foundational architectures include RNN-based approaches (GRU4Rec~\cite{hidasi2016session}), self-attentive models (SASRec~\cite{kang2018self}), and bidirectional masked transformers (BERT4Rec~\cite{sun2019bert4rec}). While collaborative models excel at capturing discrete item transition graphs, they struggle under cold-start settings and sharp behavioral transitions. Recent efforts incorporate item metadata and pre-trained language model representations~\cite{hou2022towards,petrov2023effective,wu2021multimodal}, motivating dual-expert architectures that bridge collaborative and semantic modalities.

\subsection{Mixture-of-Experts and Adaptive Gating in RecSys}
Mixture-of-Experts (MoE) architectures~\cite{shazeer2017outrageously,ma2018modeling} use learned gating networks to dynamically combine specialized sub-networks. In recommender systems, gating is widely used for multi-task learning (MMoE~\cite{ma2018modeling}, PLE~\cite{tang2020progressive}) and multi-modal feature fusion. Standard gating formulations include:
\begin{enumerate}[leftmargin=*]
    \item \textbf{Scalar Latent Gate:} $\alpha(u) = \sigma(\mathbf{w}^T [\mathbf{h}_c; \mathbf{h}_s] + b)$, directly predicting a continuous interpolation weight between two representations.
    \item \textbf{Softmax Evidence Gate:} $\mathbf{w}_u = \text{Softmax}(\mathbf{W} \mathbf{x}_u)$, allocating weights over discrete expert heads conditioned on uncertainty or entropy features~\cite{kendall2017uncertainties}.
    \item \textbf{Uncertainty-Aware Gating:} Gating conditioned on Bayesian or Monte-Carlo Dropout uncertainty estimates.
\end{enumerate}

Despite the theoretical flexibility of learned gating, empirical studies in recommender systems frequently observe that learned gating collapses toward an average blend, assigns disproportionate weight to majority patterns, or fails to outperform well-tuned static baselines. This work provides a rigorous mechanistic investigation into why this collapse occurs in dual-expert sequential settings.

\subsection{Conceptual Positioning: How CC-PoE Differs from Existing Fusion Paradigms}
\label{sec:related_positioning}

To clearly delineate our architectural positioning, Table~\ref{tab:fusion_paradigms} contrasts CC-PoE against established multi-expert and multi-modal fusion families.

\begin{table*}[t]
\centering
\small
\caption{\textbf{Conceptual Positioning of Multi-Expert Fusion Paradigms.} Contrasting classical late fusion, sequence-level gating, and candidate-level arbitration across aggregation space, disagreement handling, and theoretical guarantees.}
\label{tab:fusion_paradigms}
\begin{tabular}{lcccccc}
\toprule
\textbf{Paradigm} & \textbf{Exemplar Works} & \textbf{Level of Adaptation} & \textbf{Aggregation Space} & \textbf{Disagreement Handling} & \textbf{Theoretical Stability Guarantee} \\
\midrule
\textbf{Classical Product-of-Experts} & Hinton~\cite{hinton2002training}; Wu et al.~\cite{wu2021multimodal} & Static / Fixed ($\alpha$) & Geometric (Logit) & None (Passive Product) & Static Exponents ($\alpha > 0$) \\
\textbf{Uncertainty-Aware Gating} & Kendall \& Gal~\cite{kendall2017uncertainties} & Sequence-Level ($u$) & Arithmetic Simplex & None (Scalar Weighting) & Simplex Projection ($\sum w_k = 1$) \\
\textbf{Mixture-of-Experts (MoE / MMoE)} & Shazeer et al.~\cite{shazeer2017outrageously}; Ma et al.~\cite{ma2018modeling} & Sequence-Level ($u$) & Arithmetic / Repr. & None (Linear Convex Comb.) & Softmax Bounded $[0, 1]$ \\
\textbf{Unconstrained Item-Level MLP} & Covington et al.~\cite{covington2016deep} & Candidate-Level ($u, i$) & Unconstrained Logit & Implicit (Black-box MLP) & None (Susceptible to Tail Drift) \\
\midrule
\textbf{Adaptive CC-PoE (Ours)} & \textbf{This Work} & \textbf{Candidate-Level ($u, i$)} & \textbf{Geometric Logit} & \textbf{Piecewise Disagreement Penalty} & \textbf{\emph{A Priori} Positivity ($\min \ge 0.06 > 0$)} \\
\bottomrule
\end{tabular}
\end{table*}

\subsubsection{Algorithm Selection, Learning to Defer, and Conservative Ensembling}
Our formulation connects fundamentally to three classical machine learning frameworks:
\begin{enumerate}[leftmargin=*]
    \item \textbf{Algorithm Selection Problem:} Dating back to Rice~\cite{rice1976algorithm}, algorithm selection seeks a mapping from problem characteristics to optimal algorithms. In sequential recommendation, the sequence summary vector $\mathbf{x}_u$ serves as problem characteristics. We prove that because $\mathbf{x}_u$ lacks item-level foresight, Rice's mapping cannot be resolved at the sequence level.
    \item \textbf{Learning to Defer:} In learning to defer~\cite{madras2018predict,mozannar2020consistent}, a classifier chooses between predicting directly or deferring to a human expert. In our dual-expert setup, when sequence-level indicators are ambiguous, the system defers to the stationary geometric consensus ($0.5 \log P_C + 0.5 \log P_S$), protecting the consensus majority from misallocation.
    \item \textbf{Conservative Ensembling:} In low signal-to-noise regimes, complex gating parameters overfit to validation noise. Rather than unconstrained neural parameterization, enforcing a conservative geometric consensus inductive bias protects top-$K$ rankings against spurious tail intrusions.
\end{enumerate}
